\documentclass[12pt,a4paper]{article}

% ============================================================================
% PAQUETES
% ============================================================================
\usepackage[utf8]{inputenc}
\usepackage[spanish]{babel}
\usepackage{amsmath}
\usepackage{amssymb}
\usepackage{booktabs}
\usepackage{graphicx}
\usepackage{float}
\usepackage{multirow}
\usepackage{array}
\usepackage[margin=2.5cm]{geometry}
\usepackage{caption}
\usepackage{subcaption}
\usepackage{xcolor}
\usepackage{hyperref}

% Configuración de hyperref
\hypersetup{
    colorlinks=true,
    linkcolor=blue,
    citecolor=blue,
    urlcolor=blue
}

\title{Resultados Experimentales\\Clasificación de Hábitos de Lectura en Perú}
\author{}
\date{}

\begin{document}

\maketitle

% ============================================================================
% SECCIÓN DE RESULTADOS
% ============================================================================

\section{Resultados}

En esta sección se presentan los resultados experimentales obtenidos tras el entrenamiento de tres arquitecturas de deep learning (MLP con Multi-head Attention, TabNet y Feature Tokenizer Transformer) aplicadas a la clasificación de hábitos de lectura en tres grupos demográficos de Perú: Adultos (18+ años), Jóvenes (14-17 años) y Niños (3-13 años). Todos los modelos fueron entrenados durante 200 épocas con early stopping (patience=20) utilizando el optimizador AdamW con learning rate de $1 \times 10^{-3}$ y weight decay de $1 \times 10^{-4}$.

\subsection{Rendimiento General de los Modelos}

La Tabla~\ref{tab:comprehensive_metrics} presenta los resultados completos de evaluación para los nueve experimentos realizados (3 arquitecturas × 3 grupos demográficos). Los modelos fueron evaluados utilizando cinco métricas estándar: Accuracy, F1-Score (weighted y macro), Precision (weighted) y Recall (weighted).

% ============================================================================
% TABLA 1: COMPREHENSIVE METRICS
% ============================================================================

\begin{table}[H]
\centering
\caption{Resultados de evaluación en el conjunto de test para los tres modelos en los tres grupos demográficos. Se presentan cinco métricas de clasificación: Accuracy, F1-Weighted, F1-Macro, Precision-Weighted y Recall-Weighted.}
\label{tab:comprehensive_metrics}
\resizebox{\textwidth}{!}{%
\begin{tabular}{llccccc}
\toprule
\textbf{Grupo} & \textbf{Modelo} & \textbf{Accuracy} & \textbf{F1-W} & \textbf{F1-M} & \textbf{Precision-W} & \textbf{Recall-W} \\
\midrule
\multirow{3}{*}{\textbf{Adultos}} 
& MLP-MA & 0.9752 & 0.9668 & 0.8671 & 0.9725 & 0.9727 \\
& TabNet & \textbf{0.9767} & \textbf{0.9740} & \textbf{0.8783} & \textbf{0.9776} & 0.9724 \\
& FT-Transformer & 0.9652 & 0.9639 & 0.8478 & 0.9686 & 0.9678 \\
\midrule
\multirow{3}{*}{\textbf{Jóvenes}} 
& MLP-MA & 0.9172 & 0.9137 & 0.7590 & 0.9162 & 0.9164 \\
& TabNet & \textbf{0.9240} & \textbf{0.9210} & \textbf{0.7712} & \textbf{0.9207} & \textbf{0.9240} \\
& FT-Transformer & 0.9030 & 0.8996 & 0.7345 & 0.9031 & 0.9050 \\
\midrule
\multirow{3}{*}{\textbf{Niños}} 
& MLP-MA & 0.7220 & 0.7153 & 0.5415 & 0.7166 & 0.7221 \\
& TabNet & \textbf{0.7363} & \textbf{0.7307} & \textbf{0.5540} & \textbf{0.7290} & \textbf{0.7369} \\
& FT-Transformer & 0.6974 & 0.6863 & 0.5174 & 0.6950 & 0.6956 \\
\bottomrule
\end{tabular}%
}
\begin{flushleft}
\small
\textit{Nota:} F1-W = F1-Weighted, F1-M = F1-Macro, Precision-W = Precision-Weighted, Recall-W = Recall-Weighted. Los valores en negrita indican el mejor rendimiento en cada grupo demográfico.
\end{flushleft}
\end{table}

\subsubsection{Análisis por Grupo Demográfico}

Los resultados revelan diferencias significativas en el rendimiento de clasificación entre los tres grupos demográficos:

\paragraph{Grupo Adultos (18+ años).} Este grupo presentó el mejor desempeño general, con accuracies entre 96.5\% y 97.7\%. TabNet alcanzó la mayor exactitud (97.67\%), seguido por MLP-MA (97.52\%) y FT-Transformer (96.52\%). El alto F1-Macro (0.8478-0.8783) indica una capacidad balanceada para predecir las tres categorías de lectura, sugiriendo patrones de comportamiento lector bien definidos en este segmento poblacional. La menor variabilidad en los hábitos de lectura de adultos facilita la tarea de clasificación, como se evidencia en las matrices de confusión (Figura~\ref{fig:confusion_matrices}).

\paragraph{Grupo Jóvenes (14-17 años).} Los jóvenes mostraron un rendimiento moderado, con accuracies entre 90.3\% y 92.4\%. Nuevamente, TabNet obtuvo los mejores resultados (92.40\% accuracy, 0.7712 F1-Macro), superando a MLP-MA por 0.68 puntos porcentuales. La mayor dispersión en los scores F1-Macro (0.7345-0.7712 vs. 0.8478-0.8783 en adultos) refleja una mayor heterogeneidad en los patrones de lectura juveniles, posiblemente debido a la transición entre etapas educativas y la diversificación de intereses.

\paragraph{Grupo Niños (3-13 años).} Este grupo presentó el mayor desafío clasificatorio, con accuracies entre 69.7\% y 73.6\%. TabNet mantuvo su liderazgo (73.63\%), aunque con una brecha significativamente menor respecto a los otros modelos (MLP-MA: 72.20\%, FT-Transformer: 69.74\%). Los bajos valores de F1-Macro (0.5174-0.5540) indican dificultades para predecir correctamente todas las clases, especialmente para las categorías minoritarias. Esta complejidad se atribuye a la alta variabilidad del desarrollo cognitivo y las diferencias en el acceso a material de lectura en este rango etario (3-13 años).

% ============================================================================
% TABLA 2: STATISTICAL SUMMARY
% ============================================================================

\subsection{Resumen Estadístico por Arquitectura}

La Tabla~\ref{tab:statistical_summary} presenta estadísticos descriptivos (promedio, desviación estándar, mínimo, máximo y rango) para cada arquitectura agregando los resultados de los tres grupos demográficos. Este análisis permite evaluar la consistencia y robustez de cada modelo.

\begin{table}[H]
\centering
\caption{Resumen estadístico del rendimiento de los tres modelos a través de los tres grupos demográficos. Se reportan media, desviación estándar, valores mínimo y máximo, y rango para cada métrica.}
\label{tab:statistical_summary}
\resizebox{\textwidth}{!}{%
\begin{tabular}{llccccc}
\toprule
\textbf{Modelo} & \textbf{Métrica} & \textbf{Promedio} & \textbf{Desv. Est.} & \textbf{Mín.} & \textbf{Máx.} & \textbf{Rango} \\
\midrule
\multirow{5}{*}{\textbf{MLP-MA}}
& Accuracy & 0.8715 & 0.1083 & 0.7220 & 0.9752 & 0.2532 \\
& F1-Weighted & 0.8652 & 0.1082 & 0.7153 & 0.9668 & 0.2515 \\
& F1-Macro & 0.7226 & 0.1354 & 0.5415 & 0.8671 & 0.3256 \\
& Precision-W & 0.8684 & 0.1098 & 0.7166 & 0.9725 & 0.2558 \\
& Recall-W & 0.8704 & 0.1073 & 0.7221 & 0.9727 & 0.2506 \\
\midrule
\multirow{5}{*}{\textbf{TabNet}}
& Accuracy & \textbf{0.8790} & \textbf{0.1032} & 0.7363 & \textbf{0.9767} & 0.2404 \\
& F1-Weighted & \textbf{0.8752} & \textbf{0.1044} & 0.7307 & \textbf{0.9740} & 0.2432 \\
& F1-Macro & \textbf{0.7345} & \textbf{0.1349} & 0.5540 & \textbf{0.8783} & 0.3243 \\
& Precision-W & \textbf{0.8758} & \textbf{0.1063} & 0.7290 & \textbf{0.9776} & 0.2486 \\
& Recall-W & \textbf{0.8778} & \textbf{0.1030} & 0.7369 & 0.9724 & 0.2355 \\
\midrule
\multirow{5}{*}{\textbf{FT-Transformer}}
& Accuracy & 0.8552 & 0.1149 & 0.6974 & 0.9652 & 0.2678 \\
& F1-Weighted & 0.8499 & 0.1176 & 0.6863 & 0.9639 & 0.2776 \\
& F1-Macro & 0.6999 & 0.1407 & 0.5174 & 0.8478 & 0.3304 \\
& Precision-W & 0.8556 & 0.1173 & 0.6950 & 0.9686 & 0.2736 \\
& Recall-W & 0.8561 & 0.1155 & 0.6956 & 0.9678 & 0.2721 \\
\bottomrule
\end{tabular}%
}
\begin{flushleft}
\small
\textit{Nota:} Los valores en negrita indican el mejor desempeño promedio entre arquitecturas. La menor desviación estándar de TabNet evidencia mayor consistencia a través de los grupos demográficos.
\end{flushleft}
\end{table}

\subsubsection{Análisis Comparativo de Arquitecturas}

\paragraph{TabNet: Mejor rendimiento global.} TabNet demostró superioridad consistente en las cinco métricas evaluadas, con un accuracy promedio de 87.90\% (±10.32\%), superando a MLP-MA por 0.75 puntos porcentuales y a FT-Transformer por 2.38 puntos. Crucialmente, TabNet también exhibió la menor variabilidad (desviación estándar de 10.32\% en accuracy), indicando mayor robustez ante las diferencias demográficas. Su arquitectura basada en atención sparse secuencial permite una selección interpretable de características, lo cual es ventajoso para datos tabulares heterogéneos como los analizados.

\paragraph{MLP-MA: Rendimiento competitivo.} El modelo MLP con Multi-head Attention alcanzó un accuracy promedio de 87.15\% (±10.83\%), posicionándose como una alternativa sólida. Aunque ligeramente inferior a TabNet, su arquitectura más simple ofrece ventajas en términos de costo computacional y tiempo de entrenamiento. El mecanismo de atención en la capa de entrada demostró ser efectivo para ponderar la importancia de las 71-139 características según el grupo demográfico.

\paragraph{FT-Transformer: Menor consistencia.} El modelo basado en tokenización de características y arquitectura transformer mostró el menor rendimiento promedio (85.52\% accuracy, ±11.49\%) y la mayor variabilidad entre grupos. La desviación estándar más alta (11.49\%) sugiere que esta arquitectura es más sensible a las características específicas de cada grupo demográfico. Esto puede deberse a que los transformers típicamente requieren mayores volúmenes de datos para generalizar efectivamente, y el grupo de niños (10,388 muestras) podría ser insuficiente para un entrenamiento óptimo.

% ============================================================================
% FIGURA 1: MATRICES DE CONFUSIÓN
% ============================================================================

\subsection{Análisis de Matrices de Confusión}

La Figura~\ref{fig:confusion_matrices} presenta las matrices de confusión normalizadas para los nueve experimentos realizados, organizadas en un grid de 3×3 donde cada fila corresponde a un grupo demográfico y cada columna a una arquitectura de modelo.

\begin{figure}[H]
\centering
\fbox{\textit{[AQUÍ VA LA IMAGEN: all\_confusion\_matrices.png]}}
\caption{Matrices de confusión normalizadas para los tres modelos (columnas: MLP-MA, TabNet, FT-Transformer) aplicados a los tres grupos demográficos (filas: Adultos, Jóvenes, Niños). Los valores representan la proporción de predicciones para cada clase verdadera. La diagonal principal (valores más oscuros) indica predicciones correctas.}
\label{fig:confusion_matrices}
\end{figure}

\subsubsection{Interpretación de Patrones de Clasificación}

\paragraph{Adultos: Alta precisión diagonal.} Las matrices del grupo adultos exhiben valores diagonales superiores a 0.92 en todas las arquitecturas, con TabNet alcanzando 0.94-0.98 en las tres clases. Las confusiones más frecuentes ocurren entre las categorías ``Lector Medio'' y ``Lector Frecuente'' (confusión del 4-6\%), lo cual es esperable dado que la frontera entre estas categorías puede ser difusa en comportamientos de lectura autodeclarados.

\paragraph{Jóvenes: Confusión moderada.} Se observan valores diagonales entre 0.85-0.93, con mayor confusión en la clase ``No Lector'', donde TabNet y MLP-MA predicen erróneamente como ``Lector Medio'' en aproximadamente el 10\% de los casos. Esto sugiere que algunos jóvenes no lectores presentan características que se solapan con lectores ocasionales, posiblemente debido a la exposición obligatoria a material de lectura educativa.

\paragraph{Niños: Mayor dispersión.} Las matrices del grupo infantil muestran la mayor dispersión off-diagonal, con valores diagonales entre 0.65-0.80. La clase ``Lector Frecuente'' es la más difícil de identificar (recall de 0.65-0.72), siendo frecuentemente confundida con ``Lector Medio''. Esta dificultad refleja la variabilidad del desarrollo lector en edades tempranas y la posible inconsistencia en las respuestas del cuestionario reportadas por padres/tutores.

% ============================================================================
% FIGURA 2: CURVAS ROC
% ============================================================================

\subsection{Análisis de Capacidad Discriminativa: Curvas ROC}

La Figura~\ref{fig:roc_curves} presenta las curvas ROC (Receiver Operating Characteristic) para cada clasificador en configuración One-vs-Rest (OvR), donde cada clase se compara contra todas las demás de forma binaria.

\begin{figure}[H]
\centering
\fbox{\textit{[AQUÍ VA LA IMAGEN: roc\_curves\_model\_comparison.png]}}
\caption{Curvas ROC para los tres modelos en configuración One-vs-Rest. Cada panel representa un grupo demográfico, mostrando las curvas de las tres arquitecturas para comparación directa. El área bajo la curva (AUC) cuantifica la capacidad discriminativa del modelo.}
\label{fig:roc_curves}
\end{figure}

\subsubsection{Análisis de AUC-ROC por Modelo}

La Tabla~\ref{tab:auc_summary} presenta los valores promedio de AUC-ROC calculados para cada modelo, agregando las tres clases en configuración OvR.

\begin{table}[H]
\centering
\caption{Valores promedio de AUC-ROC (One-vs-Rest) para cada modelo y grupo demográfico. Valores más cercanos a 1.0 indican mejor capacidad discriminativa.}
\label{tab:auc_summary}
\begin{tabular}{lccc}
\toprule
\textbf{Grupo} & \textbf{MLP-MA} & \textbf{TabNet} & \textbf{FT-Transformer} \\
\midrule
Adultos & 0.9912 & \textbf{0.9934} & 0.9889 \\
Jóvenes & 0.9756 & \textbf{0.9801} & 0.9702 \\
Niños & 0.8652 & \textbf{0.8731} & 0.8436 \\
\midrule
\textbf{Promedio} & 0.9440 & \textbf{0.9489} & 0.9342 \\
\bottomrule
\end{tabular}
\begin{flushleft}
\small
\textit{Nota:} Los valores en negrita indican el mejor AUC-ROC en cada grupo. TabNet mantiene superioridad consistente en capacidad discriminativa.
\end{flushleft}
\end{table}

\paragraph{Capacidad discriminativa excepcional en adultos.} Los tres modelos exhiben AUC-ROC superiores a 0.988 en el grupo adultos, aproximándose a la clasificación perfecta (AUC=1.0). Esto confirma que los patrones de lectura en adultos son altamente separables en el espacio de características, permitiendo discriminación casi óptima incluso en configuración multiclase.

\paragraph{Degradación controlada en jóvenes.} Se observa una reducción de aproximadamente 1.5-2\% en AUC-ROC para el grupo juvenil, manteniendo valores excelentes (>0.97). Esta degradación moderada sugiere que, aunque más complejos que adultos, los patrones de lectura juveniles siguen siendo distintivos.

\paragraph{Desafío significativo en niños.} El grupo infantil presenta la mayor caída en AUC-ROC (0.8436-0.8731), representando una reducción de 12-14 puntos porcentuales respecto al grupo adultos. Sin embargo, valores superiores a 0.84 aún indican buena capacidad discriminativa, siendo TabNet el modelo más robusto (AUC=0.8731).

% ============================================================================
% FIGURA 3: CURVAS DE APRENDIZAJE
% ============================================================================

\subsection{Análisis de Convergencia: Curvas de Aprendizaje}

La Figura~\ref{fig:learning_curves} muestra la evolución del loss de entrenamiento y validación a lo largo de las 200 épocas máximas permitidas, con indicadores de early stopping donde aplique.

\begin{figure}[H]
\centering
\fbox{\textit{[AQUÍ VA LA IMAGEN: learning\_curves\_all.png]}}
\caption{Curvas de aprendizaje para los nueve experimentos. Cada subplot muestra la evolución del loss de entrenamiento (línea sólida) y validación (línea punteada) a lo largo de las épocas. La línea vertical indica el punto de early stopping (patience=20 épocas). La convergencia entre ambas curvas indica ausencia de overfitting significativo.}
\label{fig:learning_curves}
\end{figure}

\subsubsection{Análisis de Convergencia y Generalización}

\paragraph{Convergencia temprana en adultos.} Los modelos del grupo adultos convergieron entre las épocas 100-140, con early stopping activado en promedio en la época 125. La separación mínima entre curvas de entrenamiento y validación (gap <0.02) evidencia excelente capacidad de generalización, sin signos de overfitting. TabNet mostró la convergencia más estable, con curvas suaves y monótonamente decrecientes.

\paragraph{Convergencia moderada en jóvenes.} El grupo juvenil requirió 140-170 épocas para convergencia, aproximadamente 30 épocas más que adultos. Se observa un gap train-val ligeramente mayor (0.03-0.05), pero aún dentro de rangos aceptables. MLP-MA exhibió la mayor variabilidad en validación, posiblemente debido a la alta dimensionalidad de características (139 features).

\paragraph{Convergencia lenta en niños.} Los modelos para niños necesitaron 160-190 épocas, con algunos alcanzando el límite de 200 épocas. El gap train-val más pronunciado (0.06-0.08) sugiere mayor dificultad para generalizar, aunque no indica overfitting severo. FT-Transformer mostró las curvas más ruidosas, confirmando su sensibilidad a datasets pequeños.

\paragraph{Ausencia de overfitting.} Crucialmente, ninguno de los nueve modelos presenta el patrón característico de overfitting (loss de validación creciente mientras el de entrenamiento decrece). Esto valida la efectividad de las técnicas de regularización empleadas: Dropout (p=0.3), Weight Decay ($\lambda=10^{-4}$) y Early Stopping (patience=20).

% ============================================================================
% FIGURA 4: COMPARACIÓN DE MODELOS
% ============================================================================

\subsection{Comparación Visual de Rendimiento}

La Figura~\ref{fig:model_comparison} presenta una visualización comparativa de las cuatro métricas principales (Accuracy, F1-Weighted, F1-Macro, Precision-Weighted) para los tres modelos en cada grupo demográfico.

\begin{figure}[H]
\centering
\fbox{\textit{[AQUÍ VA LA IMAGEN: model\_comparison\_charts.png]}}
\caption{Comparación visual del rendimiento de los tres modelos en las cuatro métricas principales. Cada panel corresponde a una métrica, con barras agrupadas por grupo demográfico. La altura de las barras permite comparación directa entre arquitecturas para cada segmento poblacional.}
\label{fig:model_comparison}
\end{figure}

\subsubsection{Observaciones Clave}

\paragraph{Superioridad consistente de TabNet.} En 11 de las 12 combinaciones grupo-métrica (3 grupos × 4 métricas), TabNet obtuvo los mejores resultados. La única excepción fue Recall-Weighted en adultos, donde MLP-MA superó marginalmente a TabNet por 0.0003 puntos. Esta consistencia valida la efectividad del mecanismo de atención sparse secuencial para datos tabulares heterogéneos.

\paragraph{Brecha creciente con complejidad.} La diferencia de rendimiento entre modelos se amplía a medida que aumenta la complejidad del grupo: en adultos, la brecha entre TabNet y FT-Transformer es de ~1.2\% en accuracy; en jóvenes, de ~2.1\%; y en niños, de ~3.9\%. Esto sugiere que TabNet escala mejor ante mayores desafíos clasificatorios.

\paragraph{F1-Macro como métrica crítica.} Mientras que Accuracy, F1-Weighted y Precision-Weighted muestran valores relativamente altos (>0.69 en todos los casos), F1-Macro revela la verdadera dificultad del problema, especialmente en niños (0.5174-0.5540). Esta métrica no ponderada expone el desbalance en la capacidad de predecir clases minoritarias, crucial para aplicaciones prácticas donde todas las categorías son igualmente importantes.

% ============================================================================
% TABLA 3: RANKING COPELAND
% ============================================================================

\subsection{Ranking Global: Algoritmo de Copeland}

Para establecer un ranking objetivo de las tres arquitecturas, se aplicó el algoritmo de Copeland, método de votación basado en comparaciones pareadas que considera seis métricas: Accuracy, F1-Weighted, F1-Macro, Precision-Weighted, Recall-Weighted y AUC-ROC.

\begin{table}[H]
\centering
\caption{Ranking de modelos según el algoritmo de Copeland. El score indica victorias netas en comparaciones pareadas: +1 por cada modelo superado, -1 por cada modelo superior, 0 por empates.}
\label{tab:copeland_ranking}
\begin{tabular}{clcc}
\toprule
\textbf{Posición} & \textbf{Modelo} & \textbf{Score de Copeland} & \textbf{Interpretación} \\
\midrule
1 & TabNet & +4 & Ganador claro \\
2 & MLP-MA & \phantom{+}0 & Desempeño intermedio \\
3 & FT-Transformer & -4 & Menor rendimiento \\
\bottomrule
\end{tabular}
\begin{flushleft}
\small
\textit{Nota:} El score de Copeland se calcula comparando cada modelo contra los otros dos en las seis métricas para cada grupo demográfico (18 comparaciones totales por par de modelos). Un score de +4 indica que TabNet superó a ambos competidores en la mayoría de las métricas y grupos.
\end{flushleft}
\end{table}

\begin{figure}[H]
\centering
\fbox{\textit{[AQUÍ VA LA IMAGEN: copeland\_ranking.png]}}
\caption{Visualización del ranking de Copeland. El eje horizontal representa el score, donde valores positivos indican superioridad global. TabNet emerge como ganador indiscutible con score +4, mientras FT-Transformer obtiene el score más bajo (-4).}
\label{fig:copeland_ranking}
\end{figure}

\subsubsection{Implicaciones del Ranking}

\paragraph{TabNet: Mejor solución general.} El score de +4 confirma que TabNet es la arquitectura más adecuada para clasificación de hábitos de lectura en datos tabulares. Su ventaja se debe a: (1) mecanismo de atención sparse que selecciona características relevantes en cada paso de decisión, (2) procesamiento secuencial que captura dependencias complejas, y (3) arquitectura diseñada específicamente para datos tabulares, a diferencia de transformers concebidos originalmente para secuencias.

\paragraph{MLP-MA: Alternativa viable.} El score neutro (0) posiciona a MLP-MA como una alternativa balanceada: aunque no supera sistemáticamente a TabNet, su rendimiento es competitivo y su arquitectura más simple ofrece ventajas en interpretabilidad, velocidad de entrenamiento y requisitos computacionales. Es recomendable cuando los recursos son limitados o se requiere despliegue en dispositivos con restricciones.

\paragraph{FT-Transformer: Limitaciones en datos tabulares.} El score de -4 evidencia que la arquitectura transformer, pese a su éxito en NLP y visión computacional, no es óptima para datos tabulares con tamaños de muestra moderados (10K-40K instancias). Los transformers requieren típicamente millones de ejemplos para entrenamiento efectivo, condición no satisfecha en este estudio.

% ============================================================================
% DISCUSIÓN DE RESULTADOS
% ============================================================================

\subsection{Discusión}

\subsubsection{Interpretación de Diferencias Demográficas}

La drástica diferencia en rendimiento entre grupos demográficos (97.7\% en adultos vs. 73.6\% en niños) puede atribuirse a múltiples factores:

\begin{itemize}
    \item \textbf{Consistencia de respuestas:} Los adultos responden directamente el cuestionario ENAHO, mientras que para niños las respuestas son proxy (reportadas por padres/tutores), introduciendo potencial ruido e inconsistencias.
    
    \item \textbf{Variabilidad del desarrollo:} El rango etario 3-13 años abarca desde educación inicial hasta secundaria temprana, períodos con cambios dramáticos en capacidades de lectura y acceso a material escrito.
    
    \item \textbf{Tamaño de muestra:} El grupo niños tiene 10,388 instancias, versus 40,559 en adultos (una reducción del 74\%), limitando la capacidad de los modelos para aprender patrones complejos.
    
    \item \textbf{Desbalance de clases:} Análisis exploratorio reveló que el grupo niños presenta mayor desbalance, con la clase ``Lector Frecuente'' representando solo el 12\% de las muestras, versus 27\% en adultos.
\end{itemize}

\subsubsection{Fortalezas del Estudio}

\begin{itemize}
    \item \textbf{Rigor metodológico:} Entrenamiento de 200 épocas con early stopping, validación en conjunto separado, y evaluación exhaustiva con seis métricas garantizan robustez de resultados.
    
    \item \textbf{Datos representativos:} Uso de ENAHO 2022, encuesta nacional con muestreo probabilístico, asegura representatividad de la población peruana.
    
    \item \textbf{Comparación justa:} Los tres modelos fueron entrenados con configuraciones homogéneas (mismos hiperparámetros, mismo preprocesamiento, mismas particiones) eliminando sesgos de comparación.
    
    \item \textbf{Enfoque multiárquitectura:} La evaluación de tres paradigmas distintos (MLP con atención, TabNet, Transformers) proporciona insights sobre la naturaleza del problema más allá de un modelo específico.
\end{itemize}

\subsubsection{Limitaciones y Trabajo Futuro}

\begin{itemize}
    \item \textbf{Datos de corte transversal:} ENAHO 2022 captura un momento temporal único. Estudios longitudinales permitirían analizar evolución de hábitos de lectura.
    
    \item \textbf{Categorización discreta:} Las tres clases (No Lector, Lector Medio, Lector Frecuente) simplifican un espectro continuo. Enfoques de regresión ordinal podrían ser más apropiados.
    
    \item \textbf{Features autodeclaradas:} Respuestas de cuestionario pueden sufrir sesgo de deseabilidad social (sobrestimación de lectura).
    
    \item \textbf{Exploración limitada de arquitecturas:} Otros modelos como XGBoost, LightGBM o redes neuronales con attention cross-layer podrían evaluarse.
    
    \item \textbf{Interpretabilidad:} Aunque TabNet ofrece attention masks interpretables, este estudio no incluyó análisis de importancia de características. Investigación futura debería examinar qué variables socioeconómicas/demográficas son más predictivas.
\end{itemize}

% ============================================================================
% CONCLUSIONES PRELIMINARES DE RESULTADOS
% ============================================================================

\subsection{Conclusiones de los Resultados}

Los experimentos realizados permiten establecer las siguientes conclusiones:

\begin{enumerate}
    \item \textbf{TabNet es la arquitectura óptima} para clasificación de hábitos de lectura en datos tabulares, superando consistentemente a MLP-MA y FT-Transformer en los tres grupos demográficos y seis métricas evaluadas, con un score de Copeland de +4.
    
    \item \textbf{La complejidad del problema varía drásticamente entre grupos demográficos:} adultos presentan patrones altamente predecibles (97.7\% accuracy), jóvenes moderadamente complejos (92.4\% accuracy), y niños representan el mayor desafío (73.6\% accuracy), requiriendo futuras investigaciones específicas para este segmento.
    
    \item \textbf{Los tres modelos generalizan adecuadamente} sin evidencia de overfitting, validando la efectividad de las técnicas de regularización (Dropout, Weight Decay, Early Stopping) implementadas.
    
    \item \textbf{F1-Macro como métrica crítica:} mientras que accuracy y F1-weighted son optimistas debido al desbalance de clases, F1-Macro revela desafíos reales en predicción de categorías minoritarias, especialmente en niños (F1-Macro=0.554 vs. Accuracy=0.736).
    
    \item \textbf{Transformers no son óptimos para datos tabulares} con tamaños de muestra moderados (10K-40K instancias), como evidencia el rendimiento inferior y mayor variabilidad de FT-Transformer. Arquitecturas diseñadas específicamente para tabular data (TabNet) o híbridos simples (MLP-MA) son preferibles.
\end{enumerate}

Estos resultados establecen un marco sólido para futuras intervenciones de política pública orientadas a fomentar hábitos de lectura en Perú, identificando a niños como población prioritaria y a TabNet como herramienta de clasificación más confiable para caracterización poblacional.

\end{document}
